% EXAMPLE ONLY - synthetic LaTeX fixture for the LaTeX reader path.
\documentclass[11pt]{article}
\title{Carbon-Aware Scheduling of Batch Jobs in Shared GPU Clusters}
\author{Y. Zhao, P. Lindqvist}
\begin{document}
\maketitle
\begin{abstract}
Batch workloads on shared GPU clusters are scheduled for throughput, but the
electricity they consume varies by a factor of 2.1 with the hour of the day. We
present a carbon-aware scheduler that delays deferrable jobs into low-carbon
windows while holding the service-level objective. On a 96-GPU production trace
the scheduler cuts carbon emissions by 18.4\% at a throughput cost of 2.6\%.
\end{abstract}

\section{Introduction}
Shared GPU clusters now account for a large share of institutional electricity
use, yet the schedulers that decide what runs when optimise queue time only
\cite{chen2023}. Grid carbon intensity fluctuates diurnally, so the same job
started six hours later can emit 40\% less CO$_2$ without any change to the
computation itself.

\section{Related Work}
Carbon-aware scheduling has been studied for geographically distributed data
centers \cite{lin2022}, where the scheduler can migrate work between regions.
Single-cluster approaches \cite{okafor2021} reduce carbon by pausing
non-urgent work, but they report no bound on the service-level objective and
assume jobs can be restarted from a checkpoint.

\begin{figure}
\centering
\includegraphics[width=0.8\linewidth]{figures/carbon-window}
\caption{Scheduling window of the proposed carbon-aware scheduler: deferrable jobs are held in a queue until the forecast carbon intensity drops below the threshold.}
\label{fig:window}
\end{figure}

\section{Method}
We formulate the problem as a two-stage decision: a forecast stage predicts
carbon intensity over the next six hours, and a placement stage admits jobs in
order of carbon efficiency subject to a queue-time bound of 30 minutes. The
scheduler is implemented as a plug-in for an existing batch system and needs no
changes to the jobs themselves. Figure~\ref{fig:window} shows the scheduling
window and the admission threshold.

\section{Results}
Table~\ref{tab:main} reports the main comparison on a 96-GPU production trace
over 14 days. The proposed scheduler reduces emissions by 18.4\% relative to the
throughput-first baseline, at a throughput cost of 2.6\%, and it keeps the
95th-percentile queue time within the 30 minute bound.

\begin{table}
\caption{Main comparison on a 96-GPU production trace over 14 days.}
\label{tab:main}
\begin{tabular}{lccc}
Scheduler & Emissions (kg CO2) & Throughput (jobs/h) & p95 queue time (min) \\
Throughput-first baseline & 1420 & 88.5 & 22 \\
Carbon-aware (ours) & 1159 & 86.2 & 29 \\
Oracle hindsight & 1044 & 87.9 & 25 \\
\end{tabular}
\end{table}

\section{Limitations}
Our evaluation covers a single cluster and a 14 day trace, so the 18.4\%
reduction is not evidence of a general rate; the scheduler also depends on a
carbon-intensity forecast whose error is not yet characterised. Jobs that hold a
GPU for less than 5 minutes cannot be deferred, and they are scheduled
immediately.

\end{document}
